\section{Normalisation constitutive et incidences spectrales du retour}
\label{iii:sec:incidencias-espectrales}

La réponse des deux feuillets et la fonctionnelle de retour conservent deux
classes complémentaires d’information. La première comprend le produit
et le quotient des valeurs propres constitutives ; la seconde comprend les
hauteurs, les périodes et les multiplicités orientées des incidences.
Cette section réunit les opérations qui permettent de distinguer ces lectures.
Son antécédent est le transport angulaire produit par APP--TRIT--TPK,
avec ses projecteurs et sa coordinatisation relevée. Les normalisations s’appliquent
après la construction de ce transport et la conservation de son orientation.

\subsection{Séparation exacte du mode commun et du mode orienté}

Soient $r_+,r_->0$ les valeurs propres étiquetées de
\eqref{iii:eq:rchannels}. Définissons
\begin{equation}
 a=\sqrt{r_+r_-},\qquad
 \eta=\frac12\log\frac{r_+}{r_-},\qquad
 \widetilde{\mathcal V}=a^{-1}\mathcal V.
 \label{iii:eq:rec27-normalizacion}
\end{equation}
Dans la chambre $x>y>0$, on a $a\in(3/4,1)$ et $\eta>0$.
La marque $\mathcal R$ conserve $\mathcal R P_\pm=\pm P_\pm$.

\begin{proposition}[Factorisation positive avec orientation conservée]
\label{iii:prop:rec27-factorizacion}
La réponse possède la factorisation unique
\begin{equation}
 \mathcal V=a\exp(\eta\mathcal R),\qquad
 \det\widetilde{\mathcal V}=1.
 \label{iii:eq:rec27-unimodular}
\end{equation}
Ses quatre lectures constitutives satisfont
\begin{equation}
 \widehat\varepsilon=a^2e^{2\eta},\quad
 \widehat\mu=a^2e^{-2\eta},\quad
 \widehat Z=e^{-2\eta},\quad
 \widehat c=a^{-2}.
 \label{iii:eq:rec27-cuatro}
\end{equation}
Avec l’orientation fixe, l’échange de feuillets conserve $a$ et remplace
$\eta$ par $-\eta$.
\end{proposition}
\begin{proof}
Le calcul spectral donne
$\exp(\eta\mathcal R)=e^\eta P_++e^{-\eta}P_-$.
Les définitions impliquent $ae^\eta=r_+$ et $ae^{-\eta}=r_-$,
de sorte que l’on retrouve exactement $\mathcal V$.
Le produit de ses valeurs propres fixe l’unique racine positive $a$ ;
leur quotient fixe l’unique valeur réelle de $\eta$.
La substitution dans \eqref{iii:eq:four} prouve les quatre identités.
Enfin, échanger les valeurs propres inverse leur quotient et conserve
leur produit.
\end{proof}

La normalisation unimodulaire conserve toute la coordonnée orientée et
sépare la coordonnée commune. Si l’on conserve également $a$, la transformation
est réversible. Si l’on exprime uniquement $\widetilde{\mathcal V}$,
$\widehat Z$ est déterminée, tandis que $\widehat c$ nécessite le
facteur commun archivé. Cette distinction explique le rôle des feuillets
réciproques dans un lecteur de produit unitaire : un tel lecteur réalise la
composante unimodulaire. La réponse constitutive complète conserve,
en outre, son déterminant.

En particulier, les lecteurs historiques de la forme
$\mu_r^\pm=e^{\pm2u}$ et $\varepsilon_r^\pm=e^{\mp2u}$
maintiennent un produit égal à un dans chaque feuillet. Leur rapport d’impédances est
$e^{4u}$. La lecture actuelle de \eqref{iii:eq:rec27-cuatro}
maintient le produit $a^4=\widehat c^{-2}$. Le paramètre historique $u$
et la coordonnée $\eta$ ne sont identifiés que lorsqu’on déclare l’application
entre ces réalisations ; la décomposition précédente spécifie
exactement quel facteur supplémentaire permet de restituer la réponse complète.

\subsection{Fonction opératoire du moment de Catalan}

La restitution du théorème~\ref{iii:thm:restitucion-funcional-catalan}
situe la constante de Catalan à l’extrémité d’une famille fonctionnelle.
L’opération qui entre dans la réponse intérieure est
\begin{equation}
 r_\pm=
 \frac{1+s_\pm+s_\pm^2}{s_\pm(1+s_\pm)}
 \mathcal D^2\mathcal C_2(s_\pm),\qquad
 s_\pm=q_\pm^{30},\quad \mathcal D=s\frac{d}{ds}.
 \label{iii:eq:rec27-catalan-operativo}
\end{equation}
La famille et sa condition à l’origine déterminent le moment ;
l’inversion cubique détermine ses deux arguments à partir des réponses.
On conserve ainsi deux niveaux de restitution : une fonction à partir d’une fonction
et des arguments à partir d’évaluations dans une fonction déjà fixée.

La composition avec \eqref{iii:eq:rec27-normalizacion} obtient $a$ et
$\eta$ à partir de ces mêmes évaluations. La composition ultérieure avec
\eqref{iii:eq:barbero-factorization} obtient la fonctionnelle orientée
de Barbero--Immirzi. La valeur à l’extrémité $\mathcal C_2(1)$,
les réponses $\mathcal D^2\mathcal C_2(s_\pm)$ et la fonctionnelle
$\mathcal H(s_-)-\mathcal H(s_+)$ remplissent donc des fonctions
mathématiques spécifiques au sein de la même chaîne. La restitution utilise
la famille complète et les conditions d’unicité déjà démontrées.

Pour préciser quelle observation distingue les canaux, écrivons
$u_\pm=\log r_\pm$. Les variations dans une coordinatisation commune vérifient
\begin{equation}
 \begin{pmatrix}d\log\widehat c\\d\log\widehat Z\end{pmatrix}
 =\begin{pmatrix}-1&-1\\-1&1\end{pmatrix}
 \begin{pmatrix}du_+\\du_-\end{pmatrix}.
 \label{iii:eq:rec27-variaciones}
\end{equation}
Le déterminant est $-2$, de sorte que la lecture conjointe retrouve
les deux variations. Chacune séparément conserve une seule combinaison.
Cette identité exprime la fonction complémentaire de la vitesse et de
l’impédance, avec l’orientation et la coordinatisation dimensionnelle fixes.

\subsection{Trois lectures distinctes des incidences $90/120$}

Les drapeaux de \eqref{iii:eq:barbero-flags} déterminent les hauteurs
$90$ et $120$. Le défaut normalisé de leurs cardinaux est $-1/12$.
La chaîne de \eqref{iii:eq:barbero-cycle} détermine, à son tour,
les poids signés $12$ et $-1$. Chaque hauteur possède sa série de retour
$S_m(q)=q^m/(1-q^{3m})$. Conserver ces trois opérations permet de
comparer précisément les deux fonctionnelles
\begin{equation}
 L_{\rm eq}(q)=S_{90}(q)-S_{120}(q),\qquad
 L_{\rm ret}(q)=12S_{90}(q)-S_{120}(q).
 \label{iii:eq:rec27-dos-kernels}
\end{equation}
La première correspond à l’antécédent à poids égaux ; la seconde est
la lecture de la chaîne orientée actuelle.

\begin{proposition}[Normalisations du pôle de retour]
Pour $q\uparrow1$, on a
\begin{equation}
 \lim(1-q)L_{\rm eq}(q)=\frac1{1080},\qquad
 \lim(1-q)L_{\rm ret}(q)=\frac1{24},\qquad
 L_{\rm ret}-L_{\rm eq}=11S_{90}.
 \label{iii:eq:rec27-polos}
\end{equation}
\end{proposition}
\begin{proof}
La factorisation de $1-q^{3m}$ prouve
$\lim(1-q)S_m(q)=1/(3m)$.
Par linéarité, les deux coefficients sont respectivement
$1/270-1/360=1/1080$ et $12/270-1/360=1/24$.
Soustraire les deux définitions prouve la dernière égalité.
\end{proof}

Le rapport $45$ entre ces coefficients singuliers résulte des poids
spécifiés. Les fonctions conservent également leurs termes réguliers ;
le rapport des pôles n’identifie pas leurs valeurs sur tout l’intervalle.
Le défaut de cardinal, le poids de la chaîne et le coefficient singulier
sont ainsi reliés à des opérations et à des normalisations différenciées.

\clearpage
\subsection{Résolution de Fourier du motif dodécaphasique}

Sur le même cycle, le motif associé aux incidences orientées est
représenté par
\begin{equation}
 p_m=12\cos\frac{2\pi\,3m}{12}
      -\cos\frac{2\pi\,4m}{12},\qquad 0\le m<12.
 \label{iii:eq:rec27-patron}
\end{equation}
Les fréquences $3$ et $4$ réalisent les accroissements de phase $90^\circ$ et
$120^\circ$. Les arguments des fonctions trigonométriques sont
exprimés en radians. Les poids proviennent de la chaîne déjà construite ;
la représentation harmonique s’applique ensuite.

\begin{proposition}[Support et normalisation spectrale]
\label{iii:prop:rec27-fourier}
Pour la transformée non normalisée
$\widehat p_k=\sum_{m=0}^{11}p_m e^{-2\pi i km/12}$,
\begin{equation}
 \widehat p_3=\widehat p_9=72,\qquad
 \widehat p_4=\widehat p_8=-6,
 \qquad \widehat p_k=0\ \text{aux autres indices}.
 \label{iii:eq:rec27-fourier}
\end{equation}
La moyenne du motif est nulle et
$\sum_m|p_m|^2=870$.
\end{proposition}
\begin{proof}
Écrire chaque cosinus comme la demi-somme de deux exponentielles produit
les coefficients $6$ dans les caractères $3,9$ et $-1/2$ dans $4,8$.
L’orthogonalité finie
$\sum_{m=0}^{11}e^{2\pi i(j-k)m/12}=12\delta_{jk}$
donne \eqref{iii:eq:rec27-fourier}. Le coefficient d’indice zéro
s’annule. La même orthogonalité donne
$\sum|p_m|^2=(2\cdot72^2+2\cdot6^2)/12=870$.
\end{proof}

Avec la transformée divisée par $12$, les coefficients sont $6$ et
$-1/2$ ; avec la normalisation unitaire, ils sont $12\sqrt3$ et $-\sqrt3$.
Le rapport des modules $12:1$ est conservé dans les deux. La suite exacte
\begin{equation}
 2(p_0,\ldots,p_{11})
 =(22,1,-23,-2,25,1,-26,1,25,-2,-23,1)
 \label{iii:eq:rec27-patron-racional}
\end{equation}
permet de vérifier la transformée au moyen d’une arithmétique algébrique exacte.

La moyenne nulle du motif appartient à l’annulation des caractères dans le
cycle. Le défaut $-1/12$ appartient au dénombrement normalisé des incidences ;
$1/24$ appartient au pôle de la série de retour. Maintenir ces domaines
explique ce que conserve chaque lecture avant de la composer avec la paire angulaire
et avec la fonctionnelle de Barbero--Immirzi.
